\documentclass[11pt]{article}
\usepackage{fontspec}
\usepackage[french,english]{babel}
\usepackage{amsmath}
\usepackage{amssymb}
\usepackage{amsthm}
\usepackage{mathrsfs}
\usepackage{tikz-cd}
\usepackage{enumitem}

\newcommand{\printedpage}[1]{\marginpar{\footnotesize p.~#1}}
\newcommand{\msp}[1]{\marginpar{\footnotesize [#1]}}
\newcommand{\circled}[1]{\tikz[baseline=(c.base)]\node[draw,circle,inner sep=1pt](c){#1};}

\begin{document}

\printedpage{54}
% head: A. Grothendieck : Dérivateurs, Ch. XVI, Édition des Universités de Montpellier II et Paris VII
\msp{page 114}
(\textit{cf.} p. 110, case b)) and \textit{in addition} Der (3). (Thus in the case of a \textit{very} special Quillen triple, the prederivator obtained is indeed a derivator.) I will spell out the proof here.

I will denote by $(\mathcal{M},\Sigma)$ the given category of models ($\Sigma\subset\operatorname{Fl}\mathcal{M}$), by $\mathbf{D}=\mathbf{D}_{\mathcal{M},W}$ the corresponding prederivator, and by
\[ W=W_D \] % [sic] W_D -> W_{\mathbf{D}}, index printed as italic D, elsewhere bold D
the associated localizer in Cat, consisting of the ``$\mathbf{D}$-equivalences''. We see at once that $W$ is strongly saturated (true for any $W_{\mathbf{D}}$, $\mathbf{D}$ a prederivator), \textit{i.e.} satisfies Loc (1 ter), a fortiori Loc (1), on the other hand it satisfies Loc (3 bis$^*$), as a consequence of the fact that $\mathbf{D}$ satisfies Loc (4$^*$, 5$^*$). A fortiori it satisfies Loc (3$^*$), and to see that every $X$ contractible in Cat is $W$-aspheric, \textit{i.e.} $X\longrightarrow e$ in $W$, we know (\textit{cf.} §1)
\msp{page 115}
that it suffices to prove that this is so if $X$ has an initial object (axiom Loc (2$^*$) dual to axiom Loc (2)). At that point, we will know that $W$ is a basic localizer, and moreover it satisfies Loc (3$^*$ bis), hence also Loc (3 bis) (strong form of the localization axiom).

Thus \textit{we may assume} that if $S$ has an initial object, then $S$ is $W$-aspheric, \textit{i.e.} $S\longrightarrow e$ is in $W$, \textit{i.e.} $p^*:\mathbf{D}(e)\longrightarrow\mathbf{D}(S)$ is fully faithful. It amounts to the same thing to say that the adjunction arrow
\[ p_!p^*\xi\longrightarrow\xi \]
is an isomorphism, for every $\xi$ in $\mathbf{D}(e)$. Now the computation of $p_!(\eta)$ (for $\eta\in\operatorname{Ob}\mathbf{D}(S)$, \textit{i.e.} $\eta\in\operatorname{Ob}\mathcal{M}(S)$) is done as follows: we choose an arrow
\[ \eta'\xrightarrow{\Sigma_S}\eta \]
with $\eta'$ ``cofibrant'' (for the relevant cofibration structure on $\mathcal{M}(S)$), and we
\msp{page 116}
set \footnote[54]{operation $p_!$ at the level of $\mathcal{M}(S)$.}
\[ p_!(\operatorname{ho}_{W,S}(\eta))=\operatorname{ho}_W(p_!(\eta')). \]
When $\eta=p^*(\xi)$, the relevant adjunction arrow in $\operatorname{Ho}_\Sigma=\mathbf{D}(e)$ is obtained by composition
\[ p_!(\eta')\longrightarrow p_!(p^*(\xi))\xrightarrow{\text{adj.}}\xi, \tag*{$(*)$} \]
where here $p_!$ is still taken in the ordinary sense ($\mathcal{M}(S)\longrightarrow\mathcal{M}(e)=\mathcal{M}$). Now since $S$ has an initial object 0, we have
\[ \begin{aligned}
p_!(\eta') &\ \simeq\ \eta'(0)\\
p_!(p^*(\xi)) &\ \simeq\ p^*(\xi)(0)\ \simeq\ \xi,
\end{aligned} \]
and the adjunction arrow in $(*)$ is the identity, the first arrow is the arrow $\eta'(0)\longrightarrow\eta(0)$ induced by $\eta'\longrightarrow\eta$, it is in $\Sigma$, hence $(*)$ is in $\Sigma$, hence it is an isomorphism in $\operatorname{Ho}_\Sigma=\mathbf{D}(e)$, q.e.d.

The same argument works, moreover, if $\mathcal{M}$ is only a category with cofibrations or with fibrations, but then restricting oneself to the case where $S$ is a finite ordered set.
% foot: Version 18 mars 2019

% Original hard page break suppressed for the reflowed reader.
\printedpage{55}
% head: A. Grothendieck : Dérivateurs, Ch. XVI, Édition des Universités de Montpellier II et Paris VII
\msp{page 117}
I will set out separately what we have just obtained:

\textbf{Proposition 7.7.1.} \textit{Let} $(\mathcal{M},W)$ \textit{be a category of models, associated to a closed Quillen triple which is} special $(W,C,F)$\textit{. Then the prederivator} $\mathbf{D}=\mathbf{D}_{\mathcal{M},W}$ \textit{on} Cat \textit{satisfies the axioms} Der (1, 2, 3, 4$^*$, 5$^*$, 6$^*$)\textit{. If moreover the Quillen triple is} very \textit{special} (\textit{thus in particular} $C=\operatorname{Mono}(\mathcal{M})$), \textit{then it satisfies in addition} Der (4, 5, 6), i.e. \textit{it is a derivator in the full sense of the term. On the other hand, the localizer} $W=W_{\mathbf{D}}$ \textit{in} Cat \textit{associated to} $\mathbf{D}$ (\textit{in the case of a special Quillen triple and no more}) \textit{is a} basic localizer, \textit{satisfying in addition} Loc (1 ter, 3 bis).

\textbf{N.B.} What is said here for the case of a very special Quillen triple (which does not call upon the theory of basic localizers
\msp{page 118}
in Cat) shows once again that if $W$ a basic localizer in Cat satisfies Loc (6$'$ a, 8, 9), then $\mathrm{HOT}_W$ is a genuine derivator. We saw moreover in 7.6 (page 111) that we then have $W=W_{\mathrm{HOT}_W}$ \footnote[55]{But this was under the hypothesis Loc (3 bis)! Beware of the vicious circle.}, hence it follows:

\textbf{Corollary 7.7.2:} \textit{We have the implication}
\[ \text{Loc } (6'\text{ a, 8, 9})\ \Longrightarrow\ \text{Loc(1 ter, 3 bis)}\ \text{\footnotemark[56]}. \]
\footnotetext[56]{Loc (3 bis) not proved, probably false.}

Let us also note, for the record:

\textbf{Corollary 7.7.3.} \textit{Let} $(\mathcal{M},W)$ \textit{be a category of models, associated to a category with cofibrations. Then on the domain} Diag \textit{consisting of finite ordered sets, the prederivator} $\mathbf{D}=\mathbf{D}_{\mathcal{M},W}$ \textit{satisfies the axioms} Der (1, 2, 3, 4$^*$, 5$^*$, 6$^*$), \textit{and dually. Thus if} $W$ \textit{is associated both to a cofibration structure and to a fibration structure,} $\mathbf{D}$ \textit{is a derivator on} Diag.

\msp{page 119}
\texttt{[The following six lines seem out of context.]}

\textit{then the base change arrow}
\[ p^*f_!\longrightarrow f'_!q^* \]
\textit{between functors} $\mathrm{HOT}_W(X)\longrightarrow\mathrm{HOT}_W(S')$ \textit{is an isomorphism.}

This is practically trivial, considering, in the case $f$ $W$-smooth, that if $Z$ over $X$ is $W$-smooth over $X$, it is so over $S$; and in the case $p$ $W$-proper, using the last statement (``omission'') in \circled{6}, p. 199.

7.8. I finish by putting in the spotlight what, at present, seems to me the main substance of all these ``reminders'' and as it were the quintessence of the theory of basic localizers
% margin: three vertical rules printed in the left margin, beside the footnotes and down to the foot of the page
% foot: Version 18 mars 2019
\printedpage{56}
% head: A. Grothendieck : Dérivateurs, Ch. XVI, Édition des Universités de Montpellier II et Paris VII
developed up to now:
\[ \left\{ \begin{array}{rcl}
\text{Loc (3 bis)} & \Longrightarrow & \text{Der (4$^*$, 5$^*$)}\\[1ex]
\text{Loc (3 bis, 6 a)} & \Longrightarrow & \text{Der (1, 2, 3, 4$^*$, 5$^*$, 6$^*$)}\\[1ex]
\text{Loc (6$'$ a, 8, 9)} & \Longrightarrow & \text{all the axioms of derivators (+ Loc (1 ter)).}
\end{array} \right. \]

I consider the axioms Loc (3 bis), Loc (8), Loc (9) as ``harmless'' axioms - all the localizers one defines by ordinary cohomological or homotopical means, satisfy them. There remains the axiom Loc (6 a) or Loc (6$'$ a) of $W$-cofibrations (weak form),
\msp{page 120}
which visibly plays a crucial role in the theory of $\mathrm{HOT}_W$. On the other hand, the axioms Loc (4) (decomposition), Loc (5) (connectedness) and Loc (7) and its variants (and of course also Loc (6 b)) do not enter in at all! This is particularly impressive, as regards the axiom of $W$-fibrations Loc (7) and its variants (all equivalent to Loc (7) in the presence of Loc (8), with the sole exception of Loc (7 F)). As we saw, Loc (7) is not at all implied by Loc (6 a, 8, 9) ($\mathbf{W}_0$ does not satisfy it), at most it could be that Loc (7 F) is. The same is true of the connectedness axiom Loc (5), as one sees from the example of $W=\mathrm{Fl}(\mathrm{Cat})$, which satisfies all the axioms seen so far, with the sole exception
\msp{page 121}
of Loc (5). As for the question of the independence of the axiom Loc (4) from Loc (6 a, 8, 9), we are going to examine this kind of question in the following paragraph.

The set of axioms Loc (6$'$ a, 8, 9) now seems to me so important, that I feel like giving it a name, and calling a \textit{special} (basic) \textit{localizer} in Cat \footnote[57]{I hesitate whether I should not add the axiom Loc (3 bis), which ensures that $W=W_{\mathbf{D}}$, $\mathbf{D}=\mathrm{HOT}_W$.}, a basic localizer which satisfies them. These are thus also the ones for which we know that $\mathrm{HOT}_W$ is a genuine \textit{derivator.} If one cannot dismiss as ``pathological'' a structure as rich and delicate as a derivator, it follows at once that one cannot dismiss special localizers in Cat either (even if one affects to look down on the
\msp{page 122}
general basic localizers!). Indeed, $W$ is recovered uniquely with the help of $\mathbf{D}=\mathrm{HOT}_W$, or of any derivator equivalent to it, as the set of $\mathbf{D}$-equivalences in Cat \footnote[58]{But this is proved only assuming Loc (3 bis).}: it is the basic localizer (corresponding to a notion of ``universal'' equivalence, \textit{i.e.} in the category Cat of all the elementary ``spaces'') canonically associated to the derivator. It is therefore by no means a sufficiently fine fixed reference localizer, such as $\mathbf{W}_\infty$, that is fit to play a role!

This brings us again to a question of the kind already raised: if $W\subset\mathbf{W}_\infty$ is the smallest localizer which satisfies Loc (6$'$ a, 8) (or Loc (3 bis, 6$'$ a, 8)), can we say a) that $W$ is even \textit{special}, \textit{i.e.}
\msp{page 123}
satisfies Loc (9) (accessibility) as well?, and moreover b) do we have $W\neq\mathbf{W}_\infty$? I am convinced that this is the case (there is less merit for that $W$ in being different from $\mathbf{W}_\infty$,
% foot: Version 18 mars 2019
\printedpage{57}
% head: A. Grothendieck : Dérivateurs, Ch. XVI, Édition des Universités de Montpellier II et Paris VII
than for the $\tilde{\mathbf{W}}_\infty$ considered in 6.4 (p. 87)).

7.9. Omission: question 1 bis (p. 98), on the existence of a closed Quillen structure on Cat, corresponding to $W$. The only result I know in this direction is that of \textsc{Thomason}, explicitly constructing such a Quillen structure (moreover proper), in the case where $W=\mathbf{W}_\infty$. I have not really looked at the proof of \textsc{Thomason}, and do not know whether it can be extended to more general basic localizers, for example to those satisfying Loc (6$'$ a, 8, 9). If this were so, it would moreover prove that these localizers
\msp{page 124}
have enough cofibrations and enough fibrations, hence that they satisfy Loc (6 a) and Loc (7 F). (As for Loc (6 a), it seems that one can prove it directly very simply, \textit{cf.} the following paragraph. But here again, one would have to look at some examples (homological or cohomological equivalences, $\mathbf{W}_n$, \textit{i.e.} the homotopy $n$-equivalences), and see to begin with whether Loc (7 F) is satisfied.)

It does not seem to me that the theory of \textsc{Thomason} has been up to now \texttt{[illegible words]} to be very useful, technically speaking. Psychologically, I think it has been important, for giving confidence that Cat really is an excellent category of models for Hot. I was really delighted and stimulated to learn of this theorem. But the notion of a Dwyer map (which I learned in 1983 in the article of \textsc{Thomason}) has been, technically, much more important for me: it was exactly what was missing from my panoply of ``homotopical'' notions and intuitions in Cat \dots
% foot: Version 18 mars 2019

\end{document}
